\begin{table}[htbp]
\centering
\small
\caption{Six invariants for admission to the v2 series computation}
\label{tab:s11}
\begin{tabular}{>{\raggedright\arraybackslash}p{0.307\linewidth}>{\raggedright\arraybackslash}p{0.307\linewidth}>{\raggedright\arraybackslash}p{0.307\linewidth}}
\toprule
No. & Invariant & Result \\
\midrule
1 & the registry contains 1882 documents & passed \\
2 & none of the 34 excluded documents entered the matrix & passed \\
3 & raw and normalised values of the 1814 unchanged documents match v4 & passed, 114 282 values compared, zero discrepancies \\
4 & the percentile implementation reproduces v4 on the v4 composition & passed, 88 114 values \\
5 & the v5 percentiles are computed on the new reference sample & passed \\
6 & series v1 results are not overwritten & passed \\
\bottomrule
\end{tabular}
\begin{minipage}{\linewidth}\footnotesize\vspace{2pt}
\textit{Note.} The unit of analysis is the invariant. Denominator: six invariants; admission is granted only when all six pass. Data: admission manifest, revision 6. The table describes a gate; it contains no scores. Scale: none. Abbreviations: raw and normalized — the raw and normalised feature value; percentile implementation — computation of percentiles on the reference sample. Missing values: none. No correction for multiple comparisons was applied.
\end{minipage}
\end{table}
